% ============================================================================
% main.tex
% ----------------------------------------------------------------------------
% ESTE É O ARQUIVO PRINCIPAL DO MODELO. É ELE QUE VOCÊ DEVE COMPILAR.
%
% Ele apenas ORGANIZA a ordem em que as partes do documento são montadas,
% importando (\input) os demais arquivos .tex do projeto. A lógica e o
% conteúdo de cada parte ficam nos respectivos arquivos.
%
% ESTRUTURA DE PASTAS ESPERADA (relativa a este arquivo):
%   main.tex
%   edicao.tex                  <- dados de cada edição do documento
%   estrut/
%       includes.tex             <- pacotes LaTeX usados no modelo
%       config.tex                <- dados fixos da empresa e cores
%       cabecalho-rodape.tex       <- estilos de cabeçalho/rodapé
%       campos-formulario.tex      <- comandos de campos de formulário PDF
%       capa.tex                   <- layout da capa
%       folha-rosto.tex            <- layout da folha de rosto
%       sumario.tex                <- layout do sumário
%   secoes/
%       exemplo.tex                <- conteúdo de exemplo (texto comum)
%       exemplo-formulario.tex     <- conteúdo de exemplo (com formulário)
%   imag/logos/                  <- imagens dos logos usados na capa/cabeçalho
%
% VEJA O "MANUAL_DE_USO.md" PARA INSTRUÇÕES COMPLETAS DE COMO USAR ESTE MODELO.
% ============================================================================

\documentclass[12pt,a4paper]{article}

% -- Importação dos módulos do modelo ------------------------------------------
% ====================================================================
% PREÂMBULO UNIVERSAL LATEX (TEMPLATE DE INCLUDES SEGURO E COMPLETO)
% ====================================================================
% Como usar: Salve este arquivo como 'includes.tex' no mesmo diretório
% do seu documento principal e adicione \input{includes} antes de
% \begin{document}.
% ====================================================================

% --- 0. COMPILAÇÃO E FONTE ---
\usepackage[T1]{fontenc}                    % Codificação de fonte necessária para o pdfLaTeX
\usepackage[utf8]{inputenc}                 % Garante o suporte a acentuação nativa UTF-8
\usepackage{helvet}                         % Carrega a família Helvetica (equivalente visual ao Arial)
\renewcommand{\familydefault}{\sfdefault}   % Define a fonte sem serifa (Helvetica) como padrão do documento

% --- 1. IDIOMA E AJUSTES TIPOGRÁFICOS ---
\usepackage[brazil]{babel}          % Tradução de termos padrão (Figuras, Tabelas, Sumário)
\usepackage{microtype}              % Melhorias tipográficas estruturais (ajuste de espaçamento)

% --- 2. MATEMÁTICA AVANÇADA ---
\usepackage{amsmath}                    % Comandos matemáticos básicos da AMS
\usepackage{amsfonts}                   % Fontes matemáticas adicionais
\usepackage{amssymb}                    % Símbolos matemáticos diversos
\usepackage{mathtools}                  % Extensões e correções para o amsmath
\usepackage{bm}                         % Negrito para símbolos matemáticos (vetores, matrizes)
\usepackage{siunitx}                    % Formatação padronizada de unidades do SI (crucial para Engenharia)
\sisetup{output-decimal-marker = {,}}   % Configura a vírgula como separador decimal para PT-BR/PT-PT

% --- 3. LAYOUT, PÁGINA E ESPAÇAMENTO ---
\usepackage{geometry}       % Ajuste flexível das margens da página
\geometry{
    a4paper,
    top=30mm,
    bottom=20mm,
    left=30mm,
    right=20mm,
    % includehead,
    % includefoot,
    headheight=1.4cm,
    headsep=0.4cm,
    % headsep=0.6cm,
    footskip=1cm
    % footskip=1.4cm
}

\usepackage{fancyhdr}               % Customização avançada de cabeçalhos e rodapés
\usepackage{lastpage}               % Permite referenciar o número da última página (ex: "2 de 5")
\usepackage{setspace}               % Controle do espaçamento entre linhas (\onehalfspacing, etc.)
\usepackage{multicol}               % Suporte para layouts de múltiplas colunas na mesma página
\usepackage{indentfirst}            % Garante que o primeiro parágrafo de cada seção seja indentado

% --- 4. ELEMENTOS GRÁFICOS E CORES ---
\usepackage{graphicx}               % Inclusão de imagens e gráficos externos
\usepackage{transparent}            % Permite definir opacidade de imagens (ex.: logos com 70% de opacidade)
\usepackage{xcolor}                 % Suporte avançado e paletes de cores customizadas
\usepackage{tikz}                   % Ferramenta poderosa para criação de diagramas nativos
\usetikzlibrary{calc}               % Permite aritmética de coordenadas (usada na faixa do cabeçalho)
\usepackage{pgfplots}               % Geração de gráficos matemáticos de alta precisão
\pgfplotsset{compat=1.18}           % Garante compatibilidade do pgfplots

% --- 5. TABELAS PROFISSIONAIS ---
\usepackage{booktabs}               % Linhas de tabela padrão publicação (\toprule, \midrule, \bottomrule)
\usepackage{array}                  % Extensões para propriedades das colunas em tabelas
\usepackage{colortbl}               % Permite colorir linhas/colunas/células de tabelas (\rowcolor, \cellcolor)
\usepackage{multirow}               % Permite mesclar linhas em tabelas
\usepackage{tabularx}               % Tabelas com largura total fixa e colunas auto-ajustáveis (X)
\usepackage{longtable}              % Suporte para tabelas que quebram ao longo de múltiplas páginas

% --- 6. FIGURAS, LEGENDAS E FLUTUANTES ---
\usepackage{float}                  % Controle estrito de posicionamento (modificador [H])
\usepackage{caption}                % Customização estética das legendas de figuras e tabelas
\usepackage{subcaption}             % Criação de subfiguras e sub-legendas organizadas

% Legendas de figuras/tabelas usam o tamanho "menor" do padrão (Arial 10)
\DeclareCaptionFont{arialdez}{\fontsize{10}{12}\selectfont}
\captionsetup{font=arialdez, labelfont={arialdez,bf}}

% --- 7. LISTAS E ENUMERAÇÕES ---
\usepackage{enumitem}               % Customização total de listas (itemize, enumerate, description)

% --- 8. CÓDIGO-FONTE E CAIXAS DE DESTAQUE ---
\usepackage{listings}               % Inclusão e estilização de blocos de código-fonte
\usepackage{tcolorbox}              % Criação de caixas de texto coloridas e estilizadas de alta qualidade
\tcbuselibrary{skins,breakable}     % Permite caixas que quebram páginas e estilos complexos

% --- 9. TÍTULOS E SÍMBOLOS DIVERSOS ---
\usepackage{titlesec}               % Customização visual de \section, \subsection etc.
\usepackage{pifont}                 % Símbolos diversos (ex: \ding{51} para marcações de check)

% --- 10. REFERÊNCIAS E HIPERLINKS (Devem ser carregados por último) ---
\usepackage{url}                    % Formatação correta de URLs
\usepackage{hyperref}               % Transforma referências e sumários em links clicáveis interativos
\hypersetup{
    colorlinks=true,                % True para links coloridos, False para caixas coloridas
    linkcolor=black,                % Cor dos links internos (equações, seções)
    citecolor=black,                % Cor das citações bibliográficas
    urlcolor=black                  % Cor de links externos da internet
}
\usepackage{cleveref}               % Referências cruzadas inteligentes (escreve "Figura 1" automaticamente)

% Configuração de tradução do cleveref para português (garantia adicional)
\crefname{figure}{figura}{figuras}
\Crefname{figure}{Figura}{Figuras}
\crefname{table}{tabela}{tabelas}
\Crefname{table}{Tabela}{Tabelas}
\crefname{equation}{equação}{equações}
\Crefname{equation}{Equação}{Equações}

% --------------------------------------------------------------
% Aviso "Token not allowed in a PDF string (Unicode): removing
% `subscript'" do hyperref: aparece sempre que um campo de formulário
% (\TextField/\CheckBox, usados em campos-formulario.tex) é criado
% dentro de uma macro — é um comportamento interno conhecido do
% hyperref com AcroForm, inofensivo (não afeta o PDF gerado). A linha
% abaixo apenas filtra esse aviso específico da saída do compilador.
\usepackage{silence}
\WarningFilter{hyperref}{Token not allowed in a PDF string}            % pacotes LaTeX necessários (fancyhdr, hyperref, etc.)
% ============================================================================
% config.tex
% ----------------------------------------------------------------------------
% Este arquivo concentra as VARIÁVEIS DE CONFIGURAÇÃO VISUAL e os DADOS FIXOS
% da empresa/entidade que aparecem em todo o documento (capa, cabeçalho,
% rodapé, folha de rosto etc.).
%
% IMPORTANTE: este é o principal arquivo que você deve editar quando for
% adaptar o modelo para OUTRA empresa/marca (cores, logos, dados de contato).
% Os dados que mudam a cada documento específico (título, cliente, projeto)
% ficam em "edicao.tex", não aqui.
% ============================================================================

% -- Cor institucional -------------------------------------------------------
% Define a cor "verde-ej" usada nas linhas do cabeçalho/rodapé.
% Troque os valores RGB para adaptar à cor da sua empresa.
\definecolor{verde-ej}{RGB}{41, 128, 0}

% -- Caminhos dos arquivos de logo -------------------------------------------
% Devem apontar para arquivos de imagem existentes dentro do projeto.
% Se o arquivo não existir, o modelo mostra automaticamente um texto
% "(logo)" no lugar (ver uso de \IfFileExists em capa.tex e cabecalho-rodape.tex).
\newcommand{\logoCapaArquivo}{imag/logos/EletronLogo.jpg}                 % logo grande da capa
\newcommand{\logoCabecalhoArquivo}{imag/logos/EletronLogo_Horizontal.jpg} % logo pequena do cabeçalho

% -- Tamanhos dos logos -------------------------------------------------------
\newcommand{\larguraLogoCapa}{11cm}       % largura do logo na capa
\newcommand{\alturaLogoCabecalho}{1cm}  % altura do logo no cabeçalho

% -- Logos de parceiros/apoiadores (aparecem na capa, rodapé da página inicial)
\newcommand{\logoParceiraUmArquivo}{imag/logos/Logo_UnB.png}
\newcommand{\logoParceiraDoisArquivo}{imag/logos/Logo_Brasil_Junior.png}
\newcommand{\logoParceiraTresArquivo}{imag/logos/Logo_Concentro.png}

\newcommand{\alturaLogosParceiras}{1.5cm}     % altura padrão dos logos de parceiros
\newcommand{\opacidadeLogosParceiras}{0.75}    % opacidade (0 = invisível, 1 = opaco)

% -- Local e data -------------------------------------------------------------
\newcommand{\cidadeDocumento}{Brasília}  % cidade exibida na capa
\newcommand{\anoDocumento}{\the\year}    % ano atual, obtido automaticamente do sistema

% -- Dados fixos da empresa/entidade emissora ---------------------------------
% Usados no rodapé de todas as páginas do corpo do documento.
\newcommand{\razaoSocialEmpresa}{ELETRONJUN EMPRESA JÚNIOR DE ENGENHARIA ELETRÔNICA}
\newcommand{\cnpjEmpresa}{27.420.752/0001-68}
\newcommand{\empresaEndereco}{AE de Indústria Projeção A – UnB - Gama/DF - CEP: 72.444-240}
\newcommand{\telefoneEmpresa}{(61) 97402-7117}
\newcommand{\siteEmpresa}{www.eletronjun.com.br}

% Símbolo usado para separar os itens do rodapé (ex.: "Endereço | Telefone | Site")
\newcommand{\separadorRodape}{ \textbar\ }

% Nome completo da entidade executora do projeto (aparece na folha de rosto)
\newcommand{\entidadeExecutora}{EletronJun Empresa Júnior de Engenharia Eletrônica}

% Data de emissão: por padrão usa a data atual do sistema (\today).
\newcommand{\dataEmissaoDocumento}{\today}
              % dados fixos da empresa, cores e caminhos de logo
% ============================================================================
% cabecalho-rodape.tex
% ----------------------------------------------------------------------------
% Define os ESTILOS DE PÁGINA (cabeçalho e rodapé) usados no documento,
% usando o pacote "fancyhdr" (carregado em includes.tex).
%
% Este arquivo cria dois estilos de página:
%   - "corpo"   -> usado nas páginas normais de conteúdo do relatório
%   - "sumario" -> usado especificamente na página do sumário (índice)
%
% Você normalmente NÃO precisa editar este arquivo. As opções de exibição
% (mostrar CNPJ, posição da numeração) são configuradas em "edicao.tex"
% através dos comandos \NumeracaoRodapetrue/false e \CNPJnoRodapetrue/false.
% ============================================================================

% -- Linhas coloridas do cabeçalho e do rodapé --------------------------------
% Usam a cor "verde-ej" definida em config.tex.
\newcommand{\linhaCabecalhoColorida}{{\color{verde-ej}\hrule width\headwidth height 1.5pt}}
\newcommand{\linhaRodapeColorida}{{\color{verde-ej}\hrule width\headwidth height 1.5pt}}

% -- Flags (chaves liga/desliga) ----------------------------------------------
% Essas flags são definidas aqui, mas seu VALOR (true/false) é escolhido em
% edicao.tex, de acordo com o comportamento desejado para cada documento.
\newif\ifCNPJnoRodape       % controla se o CNPJ aparece no rodapé
\newif\ifNumeracaoRodape    % controla se a numeração de página fica no rodapé ou no cabeçalho

% -- Numeração de página -------------------------------------------------------
% Dependendo da flag \NumeracaoRodape, coloca o número da página no
% cabeçalho (canto superior direito) ou no rodapé (canto inferior direito).
\newcommand{\numeracaoPagina}{
    \ifNumeracaoRodape
        \fancyhead[R]{}
        \fancyfoot[R]{\fontsize{10}{12}\selectfont \raisebox{18pt}[0pt][0pt]{\thepage}}
    \else
        \fancyhead[R]{\fontsize{10}{12}\selectfont \thepage}
        \fancyfoot[R]{}
    \fi
}

% -- Texto central do rodapé ---------------------------------------------------
% Monta a linha de informações institucionais que aparece no centro do
% rodapé de cada página do corpo do documento.
% Se \CNPJnoRodape estiver ativa, inclui a razão social e o CNPJ da empresa;
% caso contrário, mostra apenas razão social e endereço.
\newcommand{\rodapeCentro}{
    \ifCNPJnoRodape
        \color{gray}\fontsize{7.5}{10}\selectfont
        \razaoSocialEmpresa\ \separadorRodape\ CNPJ: \cnpjEmpresa\ \\ \empresaEndereco \separadorRodape\ \emailEmpresa\ \separadorRodape\ \telefoneEmpresa\ \\ \siteEmpresa
    \else
        \color{gray}\fontsize{7.5}{10}\selectfont
        \razaoSocialEmpresa\ \separadorRodape\ \empresaEndereco \\ \emailEmpresa\ \separadorRodape\ \telefoneEmpresa\ \separadorRodape\ \siteEmpresa
    \fi
}

% -- Estilo de página "corpo" ---------------------------------------------------
% Usado em todas as páginas normais de conteúdo (ativado em main.tex com
% \pagestyle{corpo} logo após o sumário).
\fancypagestyle{corpo}{
    \fancyhf{}  % limpa qualquer configuração de cabeçalho/rodapé anterior

    % Linhas horizontais coloridas acima do cabeçalho e abaixo do rodapé
    \renewcommand{\headrulewidth}{1pt}
    \renewcommand{\footrulewidth}{1pt}
    \renewcommand{\headrule}{\linhaCabecalhoColorida}
    \renewcommand{\footrule}{\linhaRodapeColorida}

    % Logo no canto superior esquerdo do cabeçalho.
    % Se o arquivo de logo não existir, mostra o texto "(logo)" no lugar
    % (útil para testar o modelo antes de inserir as imagens reais).
    \fancyhead[L]{
        \IfFileExists{\logoCabecalhoArquivo}{
            \includegraphics[height=\alturaLogoCabecalho]{\logoCabecalhoArquivo}
        }{
            \small(logo)
        }
    }

    \numeracaoPagina  % posiciona a numeração conforme a flag definida em edicao.tex

    \fancyfoot[C]{\rodapeCentro}  % texto institucional no centro do rodapé
}

% -- Estilo de página "sumario" --------------------------------------------------
% Usado apenas na página do sumário/índice (\thispagestyle{sumario} em
% sumario.tex). É parecido com "corpo", mas sem numeração de página.
\fancypagestyle{sumario}{
    \fancyhf{}
    \renewcommand{\headrulewidth}{1pt}
    \renewcommand{\footrulewidth}{1pt}
    \renewcommand{\headrule}{\linhaCabecalhoColorida}
    \renewcommand{\footrule}{\linhaRodapeColorida}

    \fancyhead[L]{
        \IfFileExists{\logoCabecalhoArquivo}{
            \includegraphics[height=\alturaLogoCabecalho]{\logoCabecalhoArquivo}
        }{
            \small(logo)
        }
    }
    \fancyhead[R]{}  % sem numeração no cabeçalho da página do sumário
    \fancyfoot[R]{}  % sem numeração no rodapé da página do sumário

    \fancyfoot[C]{\rodapeCentro}
}    % define os estilos de página "corpo" e "sumario"
% ============================================================================
% edicao.tex
% ----------------------------------------------------------------------------
% Este é o arquivo que você deve editar A CADA NOVO DOCUMENTO/RELATÓRIO.
% Aqui ficam as informações que mudam de um relatório para outro:
% título, subtítulo, nome do cliente, nome do projeto, equipe, e as opções
% de exibição do rodapé.
%
% Diferente de "config.tex" (que guarda dados fixos da empresa), este
% arquivo guarda dados VARIÁVEIS de cada edição/versão do documento.
% ============================================================================

% --------------------------------------------------------------
% ONDE FICA A NUMERAÇÃO DA PÁGINA
% -------------------------------------------------------------- 
% Comente a linha que não for usar e descomente a que for utilizar.
% \NumeracaoRodapetrue  -> mostra o número da página no RODAPÉ
% \NumeracaoRodapefalse -> mostra o número da página no CABEÇALHO

% \NumeracaoRodapetrue
\NumeracaoRodapefalse


% --------------------------------------------------------------
% TEXTO NO RODAPÉ
% -------------------------------------------------------------- 
% Comente a linha que não for usar e descomente a que for utilizar.
% \CNPJnoRodapetrue  -> mostra o CNPJ da empresa no rodapé
% \CNPJnoRodapefalse -> NÃO mostra o CNPJ no rodapé

% \CNPJnoRodapetrue
\CNPJnoRodapefalse


% --------------------------------------------------------------
% EMAIL EXIBIDO NO RODAPÉ
% -------------------------------------------------------------- 
% Descomente apenas a linha do setor responsável pelo documento.

\newcommand{\emailEmpresa}{contato@eletronjun.com.br}
% \newcommand{\emailEmpresa}{projetos@eletronjun.com.br}
% \newcommand{\emailEmpresa}{marketing@eletronjun.com.br}
% \newcommand{\emailEmpresa}{gdp@eletronjun.com.br}
% \newcommand{\emailEmpresa}{operacoes@eletronjun.com.br}
% \newcommand{\emailEmpresa}{presidencia@eletronjun.com.br}
% \newcommand{\emailEmpresa}{relacionamentos@eletronjun.com.br}
% \newcommand{\emailEmpresa}{comercial@eletronjun.com.br} 
% \newcommand{\emailEmpresa}{comissaoeleitoral@eletronjun.com.br}


% --------------------------------------------------------------
% NOME DO PROJETO, CLIENTE, GERENTE E EQUIPE TÉCNICA
% -------------------------------------------------------------- 
\newcommand{\nomeProjeto}{[Nome do Projeto]}   % nome do projeto entregue
\newcommand{\nomeCliente}{[Nome do Cliente]}   % nome da empresa/cliente atendido

\newcommand{\gerenteProjeto}{[Nome do Gerente]}    % responsável pelo projeto
\newcommand{\equipeTecnica}{Assessor 1, Assessor 2} % nomes da equipe técnica (separados por vírgula)

% --------------------------------------------------------------
% TÍTULO E SUBTÍTULO DO DOCUMENTO
% -------------------------------------------------------------- 
\newcommand{\tituloDocumento}{Título do Documento}
\newcommand{\subtituloDocumento}{Subtítulo do Documento}
                     % dados específicos desta edição do documento (EDITE AQUI)
% ============================================================================
% campos-formulario.tex
% ----------------------------------------------------------------------------
% Define COMANDOS PRONTOS para criar campos de formulário PDF interativos
% (caixas de texto, checkboxes) dentro do documento, usando o pacote
% "hyperref" (que fornece \TextField e \CheckBox).
%
% Esses campos SÓ FUNCIONAM COMO INTERATIVOS quando o PDF é aberto em um
% leitor compatível com formulários (ex.: Adobe Acrobat Reader). Em alguns
% visualizadores simples (como o do navegador), os campos podem aparecer
% apenas como retângulos vazios, sem interatividade.
%
% COMANDOS DISPONÍVEIS PARA USAR NAS SEÇÕES DE CONTEÚDO:
%   \campotexto{nome}{largura}          -> campo de texto de uma linha
%   \campolongo{nome}{altura}           -> campo de texto de múltiplas linhas
%   \caixa{nome}{rótulo}                -> checkbox com texto ao lado
%   \campotabela[largura]{nome}{altura} -> campo de texto para uso dentro de tabelas
%   \checktabela{nome}                  -> checkbox para uso dentro de tabelas
%
% Exemplo de uso em uma seção de conteúdo:
%   Nome: \campotexto{nome_completo}{6cm}
%   \caixa{aceite_termos}{Li e aceito os termos}
% ============================================================================

% -- Aparência padrão dos campos -----------------------------------------------
\definecolor{fieldbg}{RGB}{245, 245, 245}   % cor de fundo cinza-clara dos campos
\newcommand{\tamanhoFonteCampo}{12pt}        % tamanho de fonte dos campos normais
\newcommand{\tamanhoFonteCampoTabela}{10pt}  % tamanho de fonte dos campos dentro de tabelas
\newcommand{\ladoCaixa}{12pt}                % tamanho (lado) dos checkboxes

% Necessário para o correto funcionamento do pacote hyperref com campos de
% formulário; não altere estas duas linhas.
\def\LayoutTextField#1#2{#2}
\def\LayoutChoiceField#1#2{#2}

% -- \campotexto{nome}{largura} ------------------------------------------------
% Cria um campo de texto de uma única linha.
%   #1 = nome interno do campo (identificador único, sem espaços)
%   #2 = largura do campo (ex.: 5cm)
\newcommand{\campotexto}[2]{
    \rule{0pt}{16pt}  % garante altura mínima de linha para alinhamento visual
    \TextField[name=#1, width=#2, height=14pt, borderwidth=0.75,
        bordercolor={0.8 0.8 0.8}, backgroundcolor=fieldbg,
        charsize=\tamanhoFonteCampo]{}
}

% -- \campolongo{nome}{altura} -------------------------------------------------
% Cria um campo de texto de múltiplas linhas (ocupa a largura toda da página).
%   #1 = nome interno do campo
%   #2 = altura do campo (ex.: 3cm)
\newcommand{\campolongo}[2]{
    \noindent\rule{0pt}{#2}
    \TextField[name=#1, width=\linewidth, height=#2, multiline=true,
        borderwidth=0.75, bordercolor={0.8 0.8 0.8}, backgroundcolor=fieldbg,
        charsize=\tamanhoFonteCampo]{}
}

% -- \caixa{nome}{rótulo} -------------------------------------------------------
% Cria um checkbox com um texto/rótulo ao lado.
%   #1 = nome interno do campo
%   #2 = texto exibido ao lado do checkbox
\newcommand{\caixa}[2]{
    {\fontsize{\tamanhoFonteCampo}{\tamanhoFonteCampo}\selectfont
    \CheckBox[name=#1, bordercolor={0 0 0},
        checkboxsymbol=\ding{51},
        width=\ladoCaixa, height=\ladoCaixa]{ #2}}
}

% -- \campotabela[largura]{nome}{altura} ---------------------------------------
% Versão do campo de texto pensada para uso DENTRO DE CÉLULAS DE TABELA.
% O parâmetro de largura é opcional (padrão: \hsize, ou seja, a largura
% disponível da célula).
%   [#1] = largura (opcional, padrão \hsize)
%   #2   = nome interno do campo
%   #3   = altura do campo
\newcommand{\campotabela}[3][\hsize]{
    \makebox[\hsize][c]{
        \TextField[name=#2, width=#1, height=#3, multiline=true,
            borderwidth=0, backgroundcolor=fieldbg,
            charsize=\tamanhoFonteCampoTabela]{}
    }
}

% -- \checktabela{nome} ---------------------------------------------------------
% Checkbox pensado para uso dentro de células de tabela (sem rótulo ao lado,
% e sem borda visível).
%   #1 = nome interno do campo
\newcommand{\checktabela}[1]{
    \makebox[\hsize][c]{
        {\fontsize{\tamanhoFonteCampoTabela}{\tamanhoFonteCampoTabela}\selectfont
        \CheckBox[name=#1, checkboxsymbol=\ding{51},
            width=\ladoCaixa, height=\ladoCaixa, borderwidth=0]{}}
    }
}   % comandos para campos de formulário PDF

\setstretch{1.15}  % espaçamento entre linhas do documento todo

% Comando auxiliar para texto em fonte menor (10pt), disponível em qualquer
% seção de conteúdo: \pequeno{texto em fonte pequena}
\newcommand{\pequeno}[1]{{\fontsize{10}{12}\selectfont #1}}

\begin{document}

% -- Páginas iniciais (sem numeração de página) --------------------------------
% ============================================================================
% capa.tex
% ----------------------------------------------------------------------------
% Gera a CAPA do documento (primeira página, sem numeração e sem cabeçalho/
% rodapé padrão).
%
% Usa os dados definidos em:
%   - config.tex  -> caminhos e tamanhos dos logos, cidade
%   - edicao.tex  -> título, subtítulo, ano
%
% Você normalmente NÃO precisa editar este arquivo, apenas os dados que ele
% utiliza (em config.tex e edicao.tex).
% ============================================================================

\begin{titlepage}
    \thispagestyle{empty}   % capa não tem cabeçalho/rodapé nem numeração
    \centering              % centraliza todo o conteúdo da capa

    % -- Logo principal da empresa --------------------------------------------
    % Se o arquivo definido em \logoCapaArquivo (config.tex) não existir,
    % mostra uma caixa com o texto "(logo da empresa aqui)" no lugar —
    % assim o modelo continua compilando mesmo sem as imagens finais.
    \IfFileExists{\logoCapaArquivo}{
        \includegraphics[width=\larguraLogoCapa]{\logoCapaArquivo}
    }{
        \fbox{\rule{0pt}{2cm}\hspace{4cm}\small (logo da empresa aqui)\hspace{4cm}}
    }

    \vspace{0.15\textheight}

    % -- Título e subtítulo do documento ---------------------------------------
    % Vêm de \tituloDocumento e \subtituloDocumento, definidos em edicao.tex.
    {\fontsize{14}{17}\selectfont\bfseries \MakeUppercase{\tituloDocumento}\par}

    \vspace{0.5cm}

    {\fontsize{12}{15}\selectfont \subtituloDocumento\par}

    \vfill

    % -- Logos de parceiros/apoiadores -----------------------------------------
    % Exibidos lado a lado, com opacidade reduzida (\opacidadeLogosParceiras).
    % Cada um também tem um "substituto" em caixa de texto, caso a imagem
    % correspondente não exista no projeto.
    \noindent
    \makebox[\textwidth]{
        \hfill
        \IfFileExists{\logoParceiraUmArquivo}{
            \transparent{\opacidadeLogosParceiras}\includegraphics[height=\alturaLogosParceiras]{\logoParceiraUmArquivo}\transparent{1}
        }{
            \fbox{\rule{0pt}{\alturaLogosParceiras}\hspace{2cm}\small(Logo UnB)\hspace{2cm}}
        }
        \hfill
        \IfFileExists{\logoParceiraDoisArquivo}{
            \transparent{\opacidadeLogosParceiras}\includegraphics[height=\alturaLogosParceiras]{\logoParceiraDoisArquivo}\transparent{1}
        }{
            \fbox{\rule{0pt}{\alturaLogosParceiras}\hspace{2cm}\small(Logo Concentro)\hspace{2cm}}
        }
        \hfill
        \IfFileExists{\logoParceiraTresArquivo}{
            \transparent{\opacidadeLogosParceiras}\includegraphics[height=\alturaLogosParceiras]{\logoParceiraTresArquivo}\transparent{1}
        }{
            \fbox{\rule{0pt}{\alturaLogosParceiras}\hspace{2cm}\small(Logo Brasil Júnior)\hspace{2cm}}
        }
        \hfill
    }

    \vspace{3cm}

    % -- Cidade e ano -----------------------------------------------------------
    {\fontsize{12}{15}\selectfont \cidadeDocumento{}, \anoDocumento\par}

\end{titlepage}

% Reinicia a contagem de páginas: a capa não é contada, então a página
% seguinte (folha de rosto) volta a ser considerada "página 1".
\setcounter{page}{1}          % capa
% ============================================================================
% folha-rosto.tex
% ----------------------------------------------------------------------------
% Gera a FOLHA DE ROSTO: a segunda página do documento, com um breve texto
% de apresentação do relatório e os dados da equipe responsável.
%
% Usa dados definidos em edicao.tex (\nomeProjeto, \nomeCliente,
% \gerenteProjeto, \equipeTecnica) e em config.tex (\entidadeExecutora,
% \dataEmissaoDocumento).
%
% Você normalmente NÃO precisa editar este arquivo, apenas os dados que ele
% utiliza (em config.tex e edicao.tex).
% ============================================================================

\thispagestyle{empty}   % folha de rosto não tem cabeçalho/rodapé

\vspace*{2cm}

% -- Texto de apresentação do relatório ---------------------------------------
\begin{center}
\begin{minipage}{0.85\textwidth}
    \setstretch{1.15}
    \noindent Relatório Técnico Conceitual apresentado como entrega
    formal do projeto \textbf{\nomeProjeto}, desenvolvido para a
    empresa \textbf{\nomeCliente}.
\end{minipage}
\end{center}

\vfill

% -- Dados da equipe responsável -----------------------------------------------
\noindent
\begin{minipage}{\textwidth}
    \textbf{Entidade Executora:} \entidadeExecutora \\[0.3cm]
    \textbf{Gerente de Projeto:} \gerenteProjeto \\[0.3cm]
    \textbf{Equipe Técnica:} \equipeTecnica
\end{minipage}

\vspace{1.5cm}

% -- Data de emissão --------------------------------------------------------
\noindent
\begin{minipage}{\textwidth}
    \textbf{Data de Emissão:} \dataEmissaoDocumento
\end{minipage}

\vspace*{1cm}

\clearpage  % garante que o próximo conteúdo (sumário) comece em página nova   % folha de rosto
% ============================================================================
% sumario.tex
% ----------------------------------------------------------------------------
% Gera a página do SUMÁRIO do documento, listando automaticamente
% as seções e subseções criadas com \section, \subsection etc. nos arquivos
% de conteúdo (pasta "secoes/").
%
% Você normalmente NÃO precisa editar este arquivo.
% ============================================================================

\pagestyle{sumario}    % estilo "sumario" (sem numeração) para toda a extensão do sumário

\tableofcontents  % gera automaticamente o sumário com base nas seções do documento
\clearpage        % começa o conteúdo principal em uma nova página

\pagestyle{corpo} % a partir daqui (conteúdo real do documento), volta ao estilo padrão
       % sumário (índice)
% Caso o documento não necessite de capa, sumário ou folha de rosto, comente as linhas acima.

% A partir daqui, todas as páginas usam o estilo "corpo": cabeçalho e
% rodapé completos, com numeração de página (ver cabecalho-rodape.tex).
\pagestyle{corpo}

% -- Conteúdo do documento ------------------------------------------------------
% Adicione aqui os arquivos de seção do seu relatório (pasta "secoes/").
% Os dois arquivos abaixo são EXEMPLOS: substitua-os pelo conteúdo real
% do seu documento, ou adicione novos \input{} para cada nova seção.
% ====================================================================
% EXEMPLO.TEX — COMO CRIAR SEÇÕES, SUBSEÇÕES E SUBSUBSEÇÕES
% ====================================================================
% Este arquivo é só um EXEMPLO de uso, não é chamado pelo main.tex.
% Para usar um destes arquivos de verdade no corpo do documento, dê
% % ====================================================================
% EXEMPLO.TEX — COMO CRIAR SEÇÕES, SUBSEÇÕES E SUBSUBSEÇÕES
% ====================================================================
% Este arquivo é só um EXEMPLO de uso, não é chamado pelo main.tex.
% Para usar um destes arquivos de verdade no corpo do documento, dê
% % ====================================================================
% EXEMPLO.TEX — COMO CRIAR SEÇÕES, SUBSEÇÕES E SUBSUBSEÇÕES
% ====================================================================
% Este arquivo é só um EXEMPLO de uso, não é chamado pelo main.tex.
% Para usar um destes arquivos de verdade no corpo do documento, dê
% \input{secoes/exemplo.tex} dentro de main.tex (onde hoje está
% \input{secoes/teste.tex}).
%
% Estrutura hierárquica disponível:
%   \section{...}       -> nível 1 (numerado 1, 2, 3...)
%   \subsection{...}    -> nível 2 (numerado 1.1, 1.2...)
%   \subsubsection{...} -> nível 3 (numerado 1.1.1, 1.1.2...)
% ====================================================================

\section{Nome da Seção}

Texto de abertura da seção. Todo o conteúdo escrito aqui, antes da
primeira \verb|\subsection|, é considerado parte da introdução desta
seção.

\subsection{Nome da Subseção}

Texto da subseção. Uma seção pode ter quantas subseções forem
necessárias, bastando repetir o comando \verb|\subsection{...}|.

\subsubsection{Nome da Subsubseção}

Texto da subsubseção — o menor nível de divisão usado neste modelo.
Uma subseção pode ter quantas subsubseções forem necessárias.

\subsubsection{Outra Subsubseção}

Basta repetir \verb|\subsubsection{...}| para criar quantas forem
necessárias dentro da mesma subseção.

\subsection{Outra Subseção}

E o mesmo vale para subseções: repita \verb|\subsection{...}| dentro
da seção quantas vezes precisar.

% --------------------------------------------------------------------
% COMO INSERIR UMA IMAGEM
% --------------------------------------------------------------------
% O comando \includegraphics (do pacote graphicx, já carregado em
% includes.tex) insere uma imagem no documento. Aqui usamos o ambiente
% "figure" para que a imagem fique centralizada e com uma legenda
% numerada automaticamente (Figura 1, Figura 2, ...).
%
%   [h]                  -> pede ao LaTeX para tentar posicionar a
%                           imagem "aqui mesmo" (here), no ponto do
%                           texto onde o comando foi escrito.
%   \centering           -> centraliza a imagem horizontalmente.
%   width=10cm           -> define a largura da imagem em centímetros
%                           (ajuste esse valor conforme o tamanho
%                           desejado; a altura é ajustada automaticamente
%                           para manter a proporção original da imagem).
%   \caption{...}        -> texto da legenda, exibido abaixo da imagem.
%   \label{fig:exemplo}  -> permite referenciar a imagem em outro ponto
%                           do texto com \ref{fig:exemplo} (ex.: "ver
%                           Figura~\ref{fig:exemplo}").
%
% O caminho da imagem é relativo à raiz do projeto (onde está o
% main.tex), então "imag/imag-exemplo.png" busca o arquivo dentro da
% pasta "imag/".
\begin{figure}[H]
    \centering
    \includegraphics[width=4cm]{imag/imag-exemplo.jpg}
    \caption{Exemplo de imagem inserida no documento.}
    \label{fig:exemplo}
\end{figure}

\section{Outra Seção}

Para começar uma nova seção (nível 1), basta usar \verb|\section{...}|
novamente. Não é necessário fechar nada antes. dentro de main.tex (onde hoje está
% \input{secoes/teste.tex}).
%
% Estrutura hierárquica disponível:
%   \section{...}       -> nível 1 (numerado 1, 2, 3...)
%   \subsection{...}    -> nível 2 (numerado 1.1, 1.2...)
%   \subsubsection{...} -> nível 3 (numerado 1.1.1, 1.1.2...)
% ====================================================================

\section{Nome da Seção}

Texto de abertura da seção. Todo o conteúdo escrito aqui, antes da
primeira \verb|\subsection|, é considerado parte da introdução desta
seção.

\subsection{Nome da Subseção}

Texto da subseção. Uma seção pode ter quantas subseções forem
necessárias, bastando repetir o comando \verb|\subsection{...}|.

\subsubsection{Nome da Subsubseção}

Texto da subsubseção — o menor nível de divisão usado neste modelo.
Uma subseção pode ter quantas subsubseções forem necessárias.

\subsubsection{Outra Subsubseção}

Basta repetir \verb|\subsubsection{...}| para criar quantas forem
necessárias dentro da mesma subseção.

\subsection{Outra Subseção}

E o mesmo vale para subseções: repita \verb|\subsection{...}| dentro
da seção quantas vezes precisar.

% --------------------------------------------------------------------
% COMO INSERIR UMA IMAGEM
% --------------------------------------------------------------------
% O comando \includegraphics (do pacote graphicx, já carregado em
% includes.tex) insere uma imagem no documento. Aqui usamos o ambiente
% "figure" para que a imagem fique centralizada e com uma legenda
% numerada automaticamente (Figura 1, Figura 2, ...).
%
%   [h]                  -> pede ao LaTeX para tentar posicionar a
%                           imagem "aqui mesmo" (here), no ponto do
%                           texto onde o comando foi escrito.
%   \centering           -> centraliza a imagem horizontalmente.
%   width=10cm           -> define a largura da imagem em centímetros
%                           (ajuste esse valor conforme o tamanho
%                           desejado; a altura é ajustada automaticamente
%                           para manter a proporção original da imagem).
%   \caption{...}        -> texto da legenda, exibido abaixo da imagem.
%   \label{fig:exemplo}  -> permite referenciar a imagem em outro ponto
%                           do texto com \ref{fig:exemplo} (ex.: "ver
%                           Figura~\ref{fig:exemplo}").
%
% O caminho da imagem é relativo à raiz do projeto (onde está o
% main.tex), então "imag/imag-exemplo.png" busca o arquivo dentro da
% pasta "imag/".
\begin{figure}[H]
    \centering
    \includegraphics[width=4cm]{imag/imag-exemplo.jpg}
    \caption{Exemplo de imagem inserida no documento.}
    \label{fig:exemplo}
\end{figure}

\section{Outra Seção}

Para começar uma nova seção (nível 1), basta usar \verb|\section{...}|
novamente. Não é necessário fechar nada antes. dentro de main.tex (onde hoje está
% \input{secoes/teste.tex}).
%
% Estrutura hierárquica disponível:
%   \section{...}       -> nível 1 (numerado 1, 2, 3...)
%   \subsection{...}    -> nível 2 (numerado 1.1, 1.2...)
%   \subsubsection{...} -> nível 3 (numerado 1.1.1, 1.1.2...)
% ====================================================================

\section{Nome da Seção}

Texto de abertura da seção. Todo o conteúdo escrito aqui, antes da
primeira \verb|\subsection|, é considerado parte da introdução desta
seção.

\subsection{Nome da Subseção}

Texto da subseção. Uma seção pode ter quantas subseções forem
necessárias, bastando repetir o comando \verb|\subsection{...}|.

\subsubsection{Nome da Subsubseção}

Texto da subsubseção — o menor nível de divisão usado neste modelo.
Uma subseção pode ter quantas subsubseções forem necessárias.

\subsubsection{Outra Subsubseção}

Basta repetir \verb|\subsubsection{...}| para criar quantas forem
necessárias dentro da mesma subseção.

\subsection{Outra Subseção}

E o mesmo vale para subseções: repita \verb|\subsection{...}| dentro
da seção quantas vezes precisar.

% --------------------------------------------------------------------
% COMO INSERIR UMA IMAGEM
% --------------------------------------------------------------------
% O comando \includegraphics (do pacote graphicx, já carregado em
% includes.tex) insere uma imagem no documento. Aqui usamos o ambiente
% "figure" para que a imagem fique centralizada e com uma legenda
% numerada automaticamente (Figura 1, Figura 2, ...).
%
%   [h]                  -> pede ao LaTeX para tentar posicionar a
%                           imagem "aqui mesmo" (here), no ponto do
%                           texto onde o comando foi escrito.
%   \centering           -> centraliza a imagem horizontalmente.
%   width=10cm           -> define a largura da imagem em centímetros
%                           (ajuste esse valor conforme o tamanho
%                           desejado; a altura é ajustada automaticamente
%                           para manter a proporção original da imagem).
%   \caption{...}        -> texto da legenda, exibido abaixo da imagem.
%   \label{fig:exemplo}  -> permite referenciar a imagem em outro ponto
%                           do texto com \ref{fig:exemplo} (ex.: "ver
%                           Figura~\ref{fig:exemplo}").
%
% O caminho da imagem é relativo à raiz do projeto (onde está o
% main.tex), então "imag/imag-exemplo.png" busca o arquivo dentro da
% pasta "imag/".
\begin{figure}[H]
    \centering
    \includegraphics[width=4cm]{imag/imag-exemplo.jpg}
    \caption{Exemplo de imagem inserida no documento.}
    \label{fig:exemplo}
\end{figure}

\section{Outra Seção}

Para começar uma nova seção (nível 1), basta usar \verb|\section{...}|
novamente. Não é necessário fechar nada antes.
\newpage
% ====================================================================
% EXEMPLO-FORMULARIO.TEX — COMO CRIAR CAMPOS, CAMPO LONGO, CAIXAS E TABELA
% ====================================================================
% Este arquivo é só um EXEMPLO de uso, não é chamado pelo main.tex.
% Para usar algo assim de verdade no corpo do documento, dê
% % ====================================================================
% EXEMPLO-FORMULARIO.TEX — COMO CRIAR CAMPOS, CAMPO LONGO, CAIXAS E TABELA
% ====================================================================
% Este arquivo é só um EXEMPLO de uso, não é chamado pelo main.tex.
% Para usar algo assim de verdade no corpo do documento, dê
% % ====================================================================
% EXEMPLO-FORMULARIO.TEX — COMO CRIAR CAMPOS, CAMPO LONGO, CAIXAS E TABELA
% ====================================================================
% Este arquivo é só um EXEMPLO de uso, não é chamado pelo main.tex.
% Para usar algo assim de verdade no corpo do documento, dê
% \input{secoes/exemplo-formulario.tex} dentro de main.tex (onde hoje
% está \input{secoes/teste.tex}).
%
% As macros usadas abaixo (\campotexto, \campolongo, \caixa,
% \campotabela, \checktabela) são definidas em
% estrut/campos-formulario.tex — não precisa mexer lá, só usá-las.
%
% REGRA OBRIGATÓRIA: todo campo só funciona dentro do ambiente Form,
% como abaixo. NOMES (1º argumento de cada macro) DEVEM SER ÚNICOS
% no documento inteiro.
% ====================================================================

\begin{Form}

\section{Exemplo de Formulário}

% --------------------------------------------------------------
% 1. Campo de texto simples (uma linha)
% --------------------------------------------------------------
Nome do responsável: \campotexto{exemplo_nome}{6cm}

\vspace{0.5cm}

% --------------------------------------------------------------
% 2. Campo de texto longo (múltiplas linhas)
% --------------------------------------------------------------
Observações gerais:

\campolongo{exemplo_observacoes}{3cm}

\vspace{0.5cm}

% --------------------------------------------------------------
% 3. Caixa de marcação (checkbox) com rótulo
% --------------------------------------------------------------
\caixa{exemplo_opcao1}{Opção 1}
\hspace{1cm}
\caixa{exemplo_opcao2}{Opção 2}
\hspace{1cm}
\caixa{exemplo_opcao3}{Opção 3}

\vspace{0.5cm}

% --------------------------------------------------------------
% 4. e 5. Campo de texto e checkbox dentro de tabela
% --------------------------------------------------------------
% Use coluna "m{...}" (em vez de "p{...}") para os campos ficarem
% centralizados também na altura da célula.
\begin{center}
\begin{tabular}{|>{\centering\arraybackslash}m{4cm}|>{\centering\arraybackslash}m{3cm}|>{\centering\arraybackslash}m{2cm}|}
    \hline
    \textbf{Item} & \textbf{Valor} & \textbf{OK?} \\
    \hline
    Item 1 & \campotabela{exemplo_tab_item1}{1cm} & \checktabela{exemplo_tab_check1} \\
    \hline
    Item 2 & \campotabela{exemplo_tab_item2}{1cm} & \checktabela{exemplo_tab_check2} \\
    \hline
\end{tabular}
\end{center}

\end{Form} dentro de main.tex (onde hoje
% está \input{secoes/teste.tex}).
%
% As macros usadas abaixo (\campotexto, \campolongo, \caixa,
% \campotabela, \checktabela) são definidas em
% estrut/campos-formulario.tex — não precisa mexer lá, só usá-las.
%
% REGRA OBRIGATÓRIA: todo campo só funciona dentro do ambiente Form,
% como abaixo. NOMES (1º argumento de cada macro) DEVEM SER ÚNICOS
% no documento inteiro.
% ====================================================================

\begin{Form}

\section{Exemplo de Formulário}

% --------------------------------------------------------------
% 1. Campo de texto simples (uma linha)
% --------------------------------------------------------------
Nome do responsável: \campotexto{exemplo_nome}{6cm}

\vspace{0.5cm}

% --------------------------------------------------------------
% 2. Campo de texto longo (múltiplas linhas)
% --------------------------------------------------------------
Observações gerais:

\campolongo{exemplo_observacoes}{3cm}

\vspace{0.5cm}

% --------------------------------------------------------------
% 3. Caixa de marcação (checkbox) com rótulo
% --------------------------------------------------------------
\caixa{exemplo_opcao1}{Opção 1}
\hspace{1cm}
\caixa{exemplo_opcao2}{Opção 2}
\hspace{1cm}
\caixa{exemplo_opcao3}{Opção 3}

\vspace{0.5cm}

% --------------------------------------------------------------
% 4. e 5. Campo de texto e checkbox dentro de tabela
% --------------------------------------------------------------
% Use coluna "m{...}" (em vez de "p{...}") para os campos ficarem
% centralizados também na altura da célula.
\begin{center}
\begin{tabular}{|>{\centering\arraybackslash}m{4cm}|>{\centering\arraybackslash}m{3cm}|>{\centering\arraybackslash}m{2cm}|}
    \hline
    \textbf{Item} & \textbf{Valor} & \textbf{OK?} \\
    \hline
    Item 1 & \campotabela{exemplo_tab_item1}{1cm} & \checktabela{exemplo_tab_check1} \\
    \hline
    Item 2 & \campotabela{exemplo_tab_item2}{1cm} & \checktabela{exemplo_tab_check2} \\
    \hline
\end{tabular}
\end{center}

\end{Form} dentro de main.tex (onde hoje
% está \input{secoes/teste.tex}).
%
% As macros usadas abaixo (\campotexto, \campolongo, \caixa,
% \campotabela, \checktabela) são definidas em
% estrut/campos-formulario.tex — não precisa mexer lá, só usá-las.
%
% REGRA OBRIGATÓRIA: todo campo só funciona dentro do ambiente Form,
% como abaixo. NOMES (1º argumento de cada macro) DEVEM SER ÚNICOS
% no documento inteiro.
% ====================================================================

\begin{Form}

\section{Exemplo de Formulário}

% --------------------------------------------------------------
% 1. Campo de texto simples (uma linha)
% --------------------------------------------------------------
Nome do responsável: \campotexto{exemplo_nome}{6cm}

\vspace{0.5cm}

% --------------------------------------------------------------
% 2. Campo de texto longo (múltiplas linhas)
% --------------------------------------------------------------
Observações gerais:

\campolongo{exemplo_observacoes}{3cm}

\vspace{0.5cm}

% --------------------------------------------------------------
% 3. Caixa de marcação (checkbox) com rótulo
% --------------------------------------------------------------
\caixa{exemplo_opcao1}{Opção 1}
\hspace{1cm}
\caixa{exemplo_opcao2}{Opção 2}
\hspace{1cm}
\caixa{exemplo_opcao3}{Opção 3}

\vspace{0.5cm}

% --------------------------------------------------------------
% 4. e 5. Campo de texto e checkbox dentro de tabela
% --------------------------------------------------------------
% Use coluna "m{...}" (em vez de "p{...}") para os campos ficarem
% centralizados também na altura da célula.
\begin{center}
\begin{tabular}{|>{\centering\arraybackslash}m{4cm}|>{\centering\arraybackslash}m{3cm}|>{\centering\arraybackslash}m{2cm}|}
    \hline
    \textbf{Item} & \textbf{Valor} & \textbf{OK?} \\
    \hline
    Item 1 & \campotabela{exemplo_tab_item1}{1cm} & \checktabela{exemplo_tab_check1} \\
    \hline
    Item 2 & \campotabela{exemplo_tab_item2}{1cm} & \checktabela{exemplo_tab_check2} \\
    \hline
\end{tabular}
\end{center}

\end{Form}

\end{document}